\documentclass[11pt,letterpaper]{article}
\usepackage[margin=0.95in,top=0.85in,bottom=0.8in]{geometry}
\usepackage[T1]{fontenc}
\usepackage{lmodern,amsmath,amssymb,amsthm,booktabs,microtype,xcolor,array,enumitem,hyperref,fancyhdr}
\hypersetup{colorlinks=true,linkcolor=teal!55!black,urlcolor=teal!55!black,citecolor=teal!55!black,pdftitle={Beyond Dirichlet GRH: a counterexample to McNew's density conjecture}}
\setlength{\headheight}{14pt}
\pagestyle{fancy}\fancyhf{}\fancyhead[L]{\small\itshape Progression-free sets and central zeros}\fancyhead[R]{\small Research note}\fancyfoot[C]{\thepage}
\newtheorem{theorem}{Theorem}[section]
\newtheorem{lemma}[theorem]{Lemma}
\theoremstyle{definition}\newtheorem{definition}[theorem]{Definition}
\newcommand{\N}{\mathbb N}\newcommand{\F}{\mathbb F}\newcommand{\Q}{\mathbb Q}
\newcommand{\grhfile}[2]{\href{https://github.com/lsernau/grh-lean/blob/72fead102f663d5c7a407964a13dc0cf9d977207/grh/#1}{#2}}
\newcommand{\incoming}[2]{\href{https://github.com/VladimirReshetnikov/ProveIt/blob/6799d49abd32ece39c01c82fc984c3a469dea648/docs/incoming/#1}{#2}}
\newcommand{\Ag}{\mathcal A}
\newcommand{\W}{\mathcal W}
\newcommand{\pr}{\mathcal P}
\setlist{nosep,topsep=4pt}
\setlength{\parskip}{1.7pt}
\raggedbottom
\begin{document}
\thispagestyle{plain}
\begin{center}
{\small\sffamily\color{teal!55!black} RESEARCH NOTE \quad / \quad 26 NOVEMBER 2026}\par
\vspace{9pt}
{\LARGE\bfseries A counterexample to McNew's\par natural-density conjecture}\par
\vspace{7pt}
{\large Carry-free prime exponents and a second transfer to central zeros}
\end{center}
\medskip
\noindent\textbf{Abstract.}
We construct a set of positive integers avoiding every nonconstant three-term geometric progression of rational ratio whose ordinary asymptotic density strictly exceeds that of Rankin's greedy set. This disproves McNew's Conjecture~7.2. A finite change to allowed base-nine digits of prime exponents gives an explicit density gain between $1.5728\times10^{-727}$ and $1.5730\times10^{-727}$. The sign also follows without computation from the divergence of the reciprocal-prime sum. The same construction improves the greedy natural density over every $\F_q[x]$ for rational-function ratios. For polynomial ratios we show that the extremal upper density and the previously conjectured value agree to every fixed order in $1/q$. Separately, using Dirichlet GRH, we prove that even with adaptive choices of $k$ real sampling sites the optimal bound within a class of positive tests for central multiplicity grows as $(\log q)/(4k)$ at fixed $k$. We also certify a zero-free central interval at one prime conductor. The density arguments use elementary arithmetic.
\section{From the Lean development to the two questions}
McNew's proposed extremizer has density about $0.71974$. Sets with larger upper density for rational ratios, and sets with larger natural density for integer ratios, are already known. His Conjecture~7.2 requires \emph{both} a natural density limit and avoidance of rational ratios \cite{mcnew}. Allowing the ratio $3/2$, for instance, forbids $4,6,9$. One prime exponent falls along this triple while another rises. Sections~2--4 give a single set that meets these requirements.

The source revision is \texttt{72fead102f66} of \texttt{lsernau/grh-lean}. Its five exported statements cover the Riemann zeta function and Dirichlet $L$-functions. They include right-half-plane nonvanishing (with the exception at the principal pole) and the critical-strip formulation for all Dirichlet characters, including imprimitive ones. The exports are in \grhfile{src/Main.lean}{\texttt{Main.lean}}, with the Dirichlet statements in \grhfile{src/Statements.lean}{\texttt{Statements.lean}}. The project uses Lean 4.35.0 and Mathlib at \texttt{56a00585}. We inspected the source and used the published independent rebuild and kernel checks for this revision \cite{sutherland}. The new proofs in this note are mathematical arguments with executable checks.

\paragraph{The combinatorial input.}
In \grhfile{src/New/Reduction/Data/Existence.lean}{\texttt{Existence.lean}, lines 77--116}, \texttt{exists\_three\_term\_free} encodes vectors as integers whose small digits preclude carries in bounded weighted midpoint relations. Restricting to a sphere forces three such vectors to coincide. The construction then uses selected exponents of a modular unit to obtain sparse sets of places (lines 167--216 and \grhfile{src/New/Reduction/Data/AdmissibleData.lean}{\texttt{IsSparse}}). The authors credit Behrend and Shapira. Our exchange uses exactly this carry-free implication, with two ternary digits grouped into one base-nine digit. The new choice is a joint set of digits across primes in which the probability gained from many new patterns exceeds the cost of removing every conflicting pattern. The density proof is independent of RH and GRH.

\paragraph{The analytic input.}
The readable edition combines short integral forms, spanning from a second moment, and a determinant lower bound to control a smoothed sum over primes. See \grhfile{GUIDE.md}{the guide}, \grhfile{src/New/Reduction/MainInequality.lean}{\texttt{star\_inequality}}, and \grhfile{src/New/Reduction/OneSidedBound.lean}{\texttt{OneSidedBound.lean}}. Its \grhfile{src/New/Reduction/Zeros/Criterion.lean}{zero criterion} exploits the signs of contributions from zero orders to logarithmic derivatives. Sections~6--7 retain those orders and concentrate a positive function on one ordinate. This addresses the additional information needed by Chowla's question about the central value. The large arithmetic-geometric estimates in the Lean proof are specific to their input. Here Dirichlet GRH enters through the location of the zeros in the Hadamard product.

\paragraph{Relation to earlier ProveIt reports.}
At revision \texttt{6799d49abd32}, the \incoming{grh_transfers_density_gaps_and_zero_tests.zip}{density-gaps report} gives a finite-prime criterion for improving the greedy natural density under rational ratios. The \incoming{mcnew_density_and_grh_transfers.zip}{McNew density report} establishes restricted optima including the four primes $2,3,5,7$, and explicitly leaves Conjecture~7.2 open. The \incoming{grh-progression-densities-and-cubic-prime-race.zip}{progression report} obtains a rational-ratio upper-density improvement over polynomial rings. The \incoming{grh-methods-polynomial-density.zip}{polynomial-density report} constructs a strict counterexample to a separate conjecture about polynomial ratios. Our construction supplies the previously missing natural-density improvement for rational ratios in both settings. Section~5 adds an upper estimate on the size of the earlier polynomial-ratio discrepancy. The \incoming{dirichlet-grh-research.zip}{critical-line report} uses an asymptotic prime-sum mean to test four central values together. Sections~6--7 instead give finite certificates for a single character and a sharp limit on their uniform strength. The comparisons identify the immediate antecedents of the present arguments.

\section{The digit exchange}
Let $\Ag=\{\sum_{j\ge0}\epsilon_j3^j:\epsilon_j\in\{0,1\},\ \epsilon_j=0
\text{ eventually}\}$. If $a+c=2b$ with $a,b,c\in\Ag$, both sides have carry-free ternary expansions. In each digit the only solutions are three zeros or three ones. Thus $a=b=c$. Rankin's set is
\[
 G=\{n\ge1:v_p(n)\in\Ag\text{ for every prime }p\}.
\]
A triple of positive integers is a nonconstant rational-ratio progression precisely when its endpoints $u,w$ and midpoint $v$ satisfy $uw=v^2$ with $u\ne w$. Unique factorization turns this into the simultaneous equalities $v_p(u)+v_p(w)=2v_p(v)$, with at least one nonconstant coordinate. This shows that $G$ is admissible. Write $d(E)=\lim_{X\to\infty}\#(E\cap[1,X])/X$ whenever the limit exists.

\begin{theorem}\label{th:main}
There exists a rational-ratio progression-free $E\subset\N$ with ordinary asymptotic density
\[
                        1.5728\times10^{-727}<d(E)-d(G)<1.5730\times10^{-727}.
\]
For every finite field $\F_q$ there is a subset of nonzero polynomials with ordinary degree-ball density strictly greater than the corresponding greedy density, avoiding every geometric progression with nonconstant ratio in $\F_q(x)$.
\end{theorem}
This interval describes the particular set constructed here. Section~3 first gives a qualitative proof and a transparent one-spoke rule. Section~4 certifies the stated numerical margin. Determining the largest possible density is a separate extremal problem.

\subsection*{A local change with an arbitrary number of primes}
Let $S$ consist of one distinguished prime and $k$ other primes. Number these coordinates $0,1,\ldots,k$. Put $C=\{0,1,3,4\}$, so that $\Ag=C+9\Ag$ uniquely. Call $C^{k+1}$ the old low-digit set. The following table specifies the new low-digit set $V$. Every unlisted coordinate in a displayed pattern has value $1$. In the one-coordinate rows only the distinguished coordinate changes.
\begin{center}
\begin{tabular}{@{}lll@{}}\toprule
Operation & Coordinates shown & Values \\ \midrule
Delete once for each value & $0$ & $1,3,4$ \\
Delete for each $i=1,\ldots,k$ & $(0,i)$ & $(0,4),(1,4),(3,4),(4,4),(4,3)$ \\
Add once & $0$ & $2$ \\
Add for each $i=1,\ldots,k$ & $(0,i)$ & $(2,2)$ \\ \bottomrule
\end{tabular}
\end{center}
The deleted patterns belonged to the old set and are all distinct. Write $x_*$ for the first added vector and $x_i$ for the added vector involving coordinate $i$.

\begin{lemma}\label{lem:digit}
$V$ contains no nonconstant vector triple $u+w=2v$, with unrestricted integral signs in $v-u$.
\end{lemma}
\begin{proof}
Three old vectors satisfy the assertion coordinate by coordinate. Two new vectors in a progression force the third vector's distinguished coordinate to be $2$, so it also is new. All new vectors lie in $\{1,2\}^{k+1}$, where distinct endpoints disagree by one in some coordinate and cannot have an integral midpoint.

It remains to exclude one new vector in a triple with two old vectors. Wherever the new digit equals $1$ and the other two digits belong to $C$, all three must equal $1$, since $C\subset\Ag$. For $x_*$ only the distinguished coordinate remains and the only surviving old pattern with every other digit $1$ has distinguished digit zero. This excludes a triple.

For $x_i$ there are two active coordinates. If $x_i$ is an endpoint, at each active coordinate the ordered old midpoint and other endpoint are either $(1,0)$ or $(3,4)$. If the other endpoint has digit $4$ at $i$, it was deleted, whatever its distinguished digit. If it has digit $0$ at $i$, the midpoint has digit $1$ there and has distinguished digit $1$ or $3$, so the midpoint was deleted. If $x_i$ is the midpoint, the old endpoint digits at $i$ are $0,4$ or $1,3$ in some order. The first possibility loses its digit-$4$ endpoint. In the second possibility the endpoint with digit $1$ at $i$ can survive only with distinguished digit $0$. The opposite endpoint then has distinguished digit $4$ and digit $3$ at $i$, exactly the final deleted pair.
\end{proof}

Define $E(S,V)$ by requiring the exponents on $S$ to have the joint form
\[
       (v_p)_{p\in S}=9h+b,\qquad h\in\Ag^{k+1},\quad b\in V .
\]
At every prime outside $S$ keep the rule $v_p\in\Ag$. If three exponent vectors are in progression, reducing their equation modulo $9$ yields $b_1+b_3=2b_2$ as an \emph{integer} vector equality since both sides lie coordinatewise between $0$ and $8$. Lemma~\ref{lem:digit} forces equality of the three low-digit vectors. Subtraction and division by nine then force equality of the high digits by progression-freeness of $\Ag$. The same conclusion already holds at the outside primes. Unique factorization proves that $E(S,V)$ is progression-free.

\section{Why the density increases}
For a prime with weight $w$ write $H(w)=\sum_{h\in\Ag}w^h$. In the integer problem $w_p=1/p$. The asymptotic probability of a finite exponent pattern $(e_p)_{p\in S}$ is $\prod_{p\in S}(1-w_p)w_p^{e_p}$, by inclusion--exclusion on multiples. Let $c=w_0$, write $w_1,\ldots,w_k$ for the remaining weights, and let $M=c\prod_i w_i$ be the unnormalized weight of the all-ones vector. Relative to $M$, altering one coordinate from $1$ to $a$ multiplies the weight by $w^{a-1}$. The gain in the sum of low-digit weights is $M B$, where
\begin{align}
 B={}&c\left(1+\sum_{i=1}^k w_i\right)-(1+c^2+c^3)\notag\\
    &-c^3\sum_{i=1}^k w_i^2
       -(c^{-1}+1+c^2+c^3)\sum_{i=1}^k w_i^3.                 \label{eq:weight}
\end{align}
The two sums on the second line pay respectively for the deleted pattern $(4,3)$ and the four patterns with digit $4$ at $i$. In the integer example take the distinguished prime to be $2$. Then
\begin{equation}\label{eq:odd}
 B=-\tfrac78+\tfrac12\sum_i p_i^{-1}
             -\tfrac18\sum_i p_i^{-2}-\tfrac{27}{8}\sum_i p_i^{-3}.
\end{equation}
The first sum diverges over the odd primes. Both negative sums converge even over all integers at least two. Therefore a finite choice makes $B>0$. This argument already disproves the conjecture without a numerical search.

\paragraph{Existence of the limit and size of the gain.}
Truncate the old rule at finitely many outside primes and cap the exponents on each selected prime at $L$. The resulting finite-coordinate conditions have densities by counting residue classes. The probability of the discarded high exponents is at most $\sum_{p\in S}p^{-L}$ (and the same argument applies to any finite list under truncation). Ignoring primes $p>Y$ outside $S$ can create an error only at multiples of $p^2$, since $0,1\in\Ag$. Uniformly in $X$, the proportion of such integers up to $X$ is bounded by $\sum_{p>Y}p^{-2}$ when normalized by $X$. Taking the two limits proves existence of the actual density. The resulting difference is exactly
\begin{equation}\label{eq:delta}
 d(E(S,V))-d(G)=M B
 \prod_{p\in S}(1-w_p)H(w_p^9)
 \prod_{p\notin S}(1-w_p)H(w_p).
\end{equation}
The outside product is positive, since each factor is at least $(1-w_p)(1+w_p)=1-w_p^2$. The selected factors are positive as well. Thus the sign of the density difference is the sign of $B$. Since natural density implies the same logarithmic density by partial summation, this construction gives the strict improvement for logarithmic density too.

\paragraph{Polynomial rings and signed rational functions.}
For a monic irreducible $P\in\F_q[x]$ take $w_P=q^{-\deg P}$. Among monic polynomials of degree $n$, the proportion divisible by $a$ is $q^{-\deg a}$ for monic $a$ with $\deg a\le n$, and is zero otherwise. Thus the same finite-dimensional limiting law holds. If $I_q(d)$ counts the degree-$d$ monic irreducibles, then $\sum_{e\mid d}eI_q(e)=q^d$, by factorization of $x^{q^d}-x$. In particular
\[
 I_q(d)\le q^d/d,\qquad
 dI_q(d)\ge q^d-\sum_{j\le d/2,\ j\ge1}q^j\ge q^d/2
 \quad\text{for all sufficiently large }d.
\]
Hence $\sum_P w_P$ diverges and $\sum_P w_P^2$ converges, for every $q$. Fixing a distinguished linear prime now makes \eqref{eq:weight} positive for an appropriate finite list of the remaining irreducibles. The same square tail proves existence of the exact-degree density, and its geometrically weighted average gives the same limit over degree balls. Formula~\eqref{eq:delta} applies with $p$ replaced by $P$.

Include all nonzero scalar multiples of each selected monic polynomial. In a progression with a rational-function ratio, the signed changes in valuations again yield vector midpoint identities. Equal valuations at every monic irreducible force the ratio to lie in $\F_q^\times$, which proves avoidance of every nonconstant rational function, including nonconstant ratios of degree zero. This strengthens the existence request in Question~6.4 of the author manuscript of Asada et al. to ordinary density \cite{asada}. In that question constant ratios are units and are excluded.

\section{Explicit certificates}
The short construction \eqref{eq:odd} already works with all $5{,}132$ odd primes through $50{,}000$. At scale $10^{15}$ the program rounds each summand in the first sum down and in the second and third sums up. It obtains the rigorous normalized bound
\[
 B\ge 543285713107/160000000000000>0.003 .
\]
The selected primorial is below $10^{21599}$ by direct integer comparison. The selected normalizing product is at least $\prod_{2\le n\le 50000}(1-1/n)=1/50000$, the outside product is at least $\prod_{n\ge2}(1-1/n^2)=1/2$, and $H(w^9)\ge1$. Thus this first certificate alone gives $d(E)-d(G)>10^{-21607}$.

In $\F_2[x]$, take $x$ as distinguished prime and all irreducibles of degrees $2$ through $16$ as remaining primes. There are $8{,}799$ selected primes including $x$, with total degree $130{,}485$. The same recurrence for $I_2(d)$ gives
\[
 B=405398067839/70368744177664>1/200 .
\]
Multiplying this normalized gain by the other factors in \eqref{eq:delta} gives the actual density difference.

\subsection*{A larger certified gain using pairs of primes}

For the integer problem a refinement gives the absolute lower bound
\[
                         d(E)-d(G)>10^{-728}.
\]
Use the same base-nine split and pivot prime $2$. Write $D=\{0,1,3,4\}$
for the old low digits. The spokes are the $265$ odd primes through
$1699$. Every new low vector has pivot $2$, with spoke digits $1$ or
$2$. Include the vector with all spoke digits $1$, all vectors with
one spoke digit $2$, and some vectors with two spoke digits $2$. The
selection of pairs is specified below.

For an old low vector, call a spoke active if its digit differs from
$1$. Delete precisely the following old vectors, with $a\in D$ the
pivot digit.
\begin{itemize}
\item With no active spoke, delete the vector if $a\in\{1,3,4\}$.
\item With one active spoke, delete its vector if the spoke digit
is $4$, or if it is $3$ and $a\ne0$.
\item With two active spokes $p<r$, delete their vector if $\{p,r\}$
is selected and their ordered digits belong to
$\{3,4\}\times\{3,4\}$ or equal $(0,4)$. In this line the selection
condition applies to both alternatives.
\end{itemize}
Keep every remaining old vector.

\paragraph{A finite proof of avoidance.}
Two new vectors in a putative progression force the third pivot to be
$2$, and the cube $\{1,2\}^{266}$ contains only constant progressions.
If there is just one new vector, every spoke where it has digit $1$ must
also be $1$ at both old points, since $D$ has no nonconstant arithmetic
progression. It suffices to check the pivot and the $k=0,1,2$ spokes
where the new vector equals $2$. For $k=2$ that pair qualifies for the
two-spoke deletion rule because its new vector was selected. Write $R$
for the test in the three bullets on these ordered coordinates, with
the pair considered selected. For each $v\in D^{k+1}$ form
\[
                  u=4\mathbf{1}-v
             \quad\text{and}\quad
                  u=2v-2\mathbf{1}.
\]
Retain $\{u,v\}$ when $u\in D^{k+1}$. The first choice treats
$2\mathbf{1}$ as the middle, the second as an endpoint with old
middle $v$. The numbers of distinct unordered pairs are respectively
$4,12,40$ for $k=0,1,2$. In every pair $R(u)$ or $R(v)$ holds.
This last finite identity is verified by the loops in
\texttt{check\_spoke\_pairs.py}. At most $4^3$ values of $v$ occur in
one case. This exhausts all signed progressions with one new point.
All-old progressions were already excluded. The base-nine lifting
preserves this conclusion since each digit sum and each doubled digit
is at most eight.

\paragraph{Selecting and pricing each pair.}
Divide the low-digit weights by the weight of the all-ones vector.
Summing over the four possible old pivot digits gives the multiplier
$2+1+1/4+1/8=27/8$. For $p<r$ the net contribution of selecting the
pair is
\[
 J(p,r)=\frac1{2pr}
 -\frac{27}{8}\left[
       (p^{-2}+p^{-3})(r^{-2}+r^{-3})+\frac{p}{r^3}\right].
\]
The last term inside the brackets is the weight of digits $(0,4)$.
Select exactly the pairs with $J(p,r)>0$. Their additions and
deletions have distinct active spoke sets, so their contributions add
without overlaps. The complete normalized gain is
\[
 B=-\frac78+
 \sum_{3\le p\le1699}\left(\frac1{2p}
                   -\frac{11}{8p^2}-\frac{27}{8p^3}\right)
             +\sum_{3\le p<r\le1699}\max\{0,J(p,r)\},
\]
where the indices in both sums are primes. Here $11/8$ is the sum of
the normalized weights of pivot digits $1,3,4$, which are the three
pivot values deleted when the single active spoke has digit $3$.

The integer checker finds $15{,}693$ selected pairs. It bounds $B$ from
below by rounding every summand down to a multiple of $10^{-15}$ and
obtains
\[
                          B\ge\frac{73099512201}{10^{15}}.
\]
For independent reproduction,
$8p^3r^3 J(p,r)=4p^2r^2-27((p+1)(r+1)+p^4)$. Put $P=\prod_{p\le1699}p$. The low-block weight
normalization including the geometric prefactors is $\varphi(P)/P^2$.
The independent higher-digit weights contribute
$\prod_{p\mid P}\sum_{h\in\Ag}p^{-9h}\ge1$. The outside greedy
density is at least $1/\zeta(2)>1/2$. Finally the checker verifies
the integer comparison
\[
       73099512201\,\varphi(P)\,10^{728}>2\cdot10^{15} P^2.
\]
Multiplying these bounds gives the claimed improvement in natural
density. No floating-point arithmetic enters the certificate.

\paragraph{The actual size of the exchange.}
Let $\eta=(d(E)-d(G))/d(G)$. The identity
$H(w)=(1+w)(1+w^3)H(w^9)$ cancels both the outside factors and all high-digit factors when dividing \eqref{eq:delta} by $d(G)$. Thus the \emph{relative} increase is an exact rational number,
\[
 \eta=B\prod_{p\le1699}\frac{p^3}{(p+1)(p^3+1)} .
\]
Computing $B$ from the displayed finite sum as a rational gives
$2.185320021572<10^{727}\eta<2.185320021573$.
We also certify
\begin{equation}\label{eq:numericgain}
              1.5728 <10^{727}\bigl(d(E)-d(G)\bigr)< 1.5730 .
\end{equation}
For the absolute interval, take $T=100000$ and multiply the factors
$(1-p^{-1})(1+p^{-1})(1+p^{-3})(1+p^{-9})$ for $p\le T$.
This retains the first three ternary positions. Relative to this truncated probability, the true $d(G)$ can be lower by at most
$\sum_{n>T}n^{-2}\le1/T$, or higher by at most
$\sum_{n\ge2}n^{-27}\le 2^{-27}+2^{-26}/26$.
The checker bounds the truncated probability by rounding each product outwards to $18$ decimal places. It tests both ends of \eqref{eq:numericgain} by rational comparisons. This completes Theorem~\ref{th:main}.

For an example of the actual rule, let $P$ be the primorial through $1699$. The exchange drops $P$ and admits $2P$, $2pP$ for each spoke, and $2prP$ for each selected pair. The high digits and primes outside the finite list supply all the other integers represented by these patterns. The change in membership therefore has a genuine limiting density even though the numerical gain is exceptionally small.


\section{The polynomial-ratio supremum at large field size}
\label{sec:charging}
Let $\beta_q$ be the supremum of upper densities of subsets of
$\mathbb F_q[x]$ containing no triple $f,fr,fr^2$ with $f\ne0$ and
$\deg r\ge1$. Densities use balls of bounded degree. Write $A$ for the
nonnegative integers with ternary digits in $\{0,1\}$. The degree-set
construction of Asada, Fourakis, Manski, McNew, Miller and Moreland
(Theorem~4.3 and Corollary~4.4 of the author manuscript,
\href{https://doi.org/10.1016/j.ffa.2016.10.002}{Finite Fields Appl.\ 44
(2017), 135--147}) gives
\[
 m_q=(1-q^{-1})\sum_{j\in A}q^{-j}\ \le\ \beta_q .
\]
The earlier ProveIt counterexample proves strict inequality. The next
estimate locates the size of the discrepancy. Its proof uses only polynomial
factorization and counting.

\paragraph{Theorem.}
For every prime power $q\ge16$,
\[
 0\le \beta_q-m_q\le (2q+1)q^{-q+14}.
\]
In particular $\beta_q$ and $m_q$ agree to every fixed order in $1/q$ as
$q$ tends to infinity through prime powers. The upper limit of
$\log(\beta_q-m_q)/(q\log q)$ is at most $-1$.

\paragraph{Proof.}
Put $K=q$ and fix a degree cap $N\ge K$ and a progression-free set $E$.
Call degree $N-k$ layer $k$, and call a layer preferred if $k\in A$.
For each $k\in[0,K]\setminus A$, form $a,b$ by replacing every digit
$2$ of $k$ by $0,1$, respectively, keeping the other digits.
Then $a,b\in A$, $a<b<k$, and $k+a=2b$.
Set $d=k-b$. For every $f\in E$ on layer $k$ and every monic
irreducible $P$ of degree $d$, at least one of $fP$ on layer $b$ and
$fP^2$ on layer $a$ is missing from $E$. Assign weight $1/I_d$ to one
such missing polynomial, where $I_d$ counts the monic irreducibles of
degree $d$. The weights sent by each $f$ sum to one. All three polynomials
have the same leading coefficient, so counting also applies to sets which
distinguish associates.

For a nonzero polynomial $g$, let $\omega_d(g)$ count its distinct monic
irreducible factors of degree $d$. For fixed $k$ and target $g$, at most
$\omega_d(g)$ choices of $P$ can send weight to $g$. Its layer determines
whether to recover $f$ as $g/P$ or $g/P^2$.
In fact a target layer $h$ and a ratio degree $d$ determine at most
one source layer $k$. The ternary digits of $d$ are ones precisely
where $k$ has twos. On this nonempty set of positions $b$ has only ones
and $a$ has only zeros. Thus $h$ cannot occur in both roles for the same
$d$. The role fixes $k$ as $h+d$ or $h+2d$.

Choose
\[
 \theta_1=1-8/q,\qquad \theta_d=\frac{8}{q d(d-1)}
 \quad(d\ge2),\qquad \sum_{d\ge1}\theta_d=1.
\]
Call $(k,g)$ exceptional if $\omega_d(g)>\theta_d I_d$.
Across all its nonexceptional pairs, $g$ receives at most one unit
of weight, by uniqueness of the source layer for each $d$.
Hence the missing elements on preferred layers can account for all
nonexceptional weight. The weight at an exceptional pair is at most one.

For $g$ uniform on any prescribed degree $n$ (also with leading
coefficient fixed), the probability that each of $t$ prescribed distinct
degree-$d$ primes divides $g$ is at most $q^{-dt}$. It equals that value
if $dt\le n$ and is zero otherwise. Taking
$t=\lfloor\theta_d I_d\rfloor+1$ and using the union bound gives
\[
 \Pr\big(\omega_d(g)\ge t\big)\le {I_d\choose t}q^{-dt}.
\]
For $d=1$ we have $I_1=q$ and $t=q-7$. The probability is at most
$\binom q7q^{-(q-7)}\le q^{-q+14}$. For $d\ge2$,
$dI_d\ge q^d-\sum_{j=1}^{\lfloor d/2\rfloor}q^j\ge q^d/2$ since
every element of $\mathbb F_{q^d}$ with degree less than $d$ belongs
to a proper subfield of degree at most $d/2$.
Hence
\[
  t>\frac{4q^{d-1}}{d^2(d-1)}\ge q,\qquad
  \frac{\mathrm e}{\theta_d q^d}
    =\frac{\mathrm e\,d(d-1)}{8q^{d-1}}\le\frac{\mathrm e}{4q}<\frac1q .
\]
The two elementary bounds hold first at $d=2$. When $d$ increases by
one, the denominators $d^2(d-1)$ and $d(d-1)$ grow by factors at
most $9/2$ and $3$, respectively, while $q\ge16$.
Using $\binom{I_d}{t}\le(\mathrm e I_d/t)^t$ now bounds the probability
by $q^{-q}$ for $d\ge2$.

Write $B_N$ for the nonzero polynomials of degree at most $N$.
Across all $k\le K$ the exceptional weight is at most
$2q|B_N|q^{-q+14}$. The proportion of $B_N$ on layers deeper than
$K$ is $(q^{N-K}-1)/(q^{N+1}-1)\le q^{-K-1}$.
The count of $E$ is at most the full count on preferred layers,
plus the exceptional weight and this deeper-layer allowance. Every
remaining member on a nonpreferred layer has one unit of weight, which
has already been charged against a missing preferred-layer element.
The proportion of $B_N$ on preferred layers tends to
$(1-q^{-1})\sum_{j\in A}q^{-j}=m_q$. Taking the upper limit in $N$,
then the supremum over $E$, and using $K=q$ proves the bound.
\hfill$\square$

\section{A finite test for Chowla's central nonvanishing question}

Let $\chi$ be a primitive nonprincipal Dirichlet character of conductor
$q$, put $b=(1-\chi(-1))/2$, and write
\[
 \Lambda_\chi(s)=(q/\pi)^{(s+b)/2}\Gamma((s+b)/2)L(s,\chi).
\]
Chowla's question in the quadratic case asks whether every such real
character satisfies $L(1/2,\chi)\ne0$. The repository's one-sided
remainder estimate permits central zeros. Its use of logarithmic derivatives
to convert arithmetic information into zero multiplicities suggests keeping
the residues themselves and using a positive kernel concentrated at a
chosen ordinate. The following kernel uses only values in the half-plane
of absolute convergence. Throughout this section we use Dirichlet GRH.

\paragraph{The finite criterion.}
Fix a real $t$, let $m_t=\operatorname{ord}_{1/2+it}L(s,\chi)$, and put
\[
 H_j(t)=\sqrt j\,\Re\frac{\Lambda_\chi'}{\Lambda_\chi}
       (1/2+\sqrt j+it),\qquad
 B_k(t)=\sum_{j=1}^k(-1)^{j-1}\binom{k}{j} H_j(t),\qquad
 K_k(u)=\prod_{j=1}^k\frac{j}{j+u}.
\]
For every $k\ge1$, counting ordinates $\gamma$ with multiplicity, we have
\begin{equation*}
 B_k(t)=\sum_\gamma K_k((\gamma-t)^2)\ \downarrow\ m_t.
\end{equation*}
Consequently $m_t\le r$ if and only if $B_k(t)<r+1$ for some finite $k$,
for each integer $r\ge0$. For a real quadratic character, the functional
equation has sign $+1$, so $m_0$ is even. In that case $B_k(0)<2$ certifies
central nonvanishing. More generally $B_k(0)<2K_k(g^2)$ for $g\ge0$
certifies the absence of all nontrivial zeros with $|\gamma|\le g$.

\emph{Proof.}
The order-one Hadamard product gives the logarithmic derivative as a
constant plus $\sum_\rho(1/(s-\rho)+1/\rho)$. The sum of
$\Re(1/\rho)$ converges absolutely, by the usual zero count. On the
critical line, away from the zeros, the functional equation
$\Lambda_\chi(s)=\epsilon_\chi\overline{\Lambda_\chi(1-\bar s)}$
makes $\Lambda_\chi'/\Lambda_\chi$ purely imaginary. Since every
$\rho=1/2+i\gamma$, this identifies the real constant in the product and gives
\[
 H_j(t)=\sum_\gamma\frac{j}{j+(\gamma-t)^2}.
\]
Partial fractions give
$\sum_{j=1}^k(-1)^{j-1}\binom{k}{j}j/(j+u)=K_k(u)$.
For $u>0$ the product tends to zero by divergence of $\sum 1/j$.
Dominated convergence applies since the sum for $H_1(t)$ is finite.
Nonreal zeros of the real character occur in pairs of opposite ordinates.
A pair with $|\gamma|\le g$, or a central zero of even order, contributes
at least $2K_k(g^2)$. This proves the claims.

\paragraph{What the convergence rate measures.}
Let $h$ be the least positive distance $|\gamma-t|$ and let $d_h$ count
the zeros at that distance. Then
\[
  B_k(t)-m_t\sim d_h\,\Gamma(1+h^2)k^{-h^2}.
\]
Indeed $K_k(u)=\Gamma(k+1)\Gamma(1+u)/\Gamma(k+1+u)$.
For $u\ge h^2$ the quotient $K_k(u)/K_k(h^2)$ is decreasing with $k$,
is bounded by $K_1(u)/K_1(h^2)$, and tends to zero for $u>h^2$.
The bound is summable over the zeros, so division by $K_k(h^2)$ and
dominated convergence give the result. A zero very close to the target
ordinate therefore slows the convergence by a precise exponent.

\paragraph{The conductor term for a fixed number of samples.}
At $t=0$ and fixed $k$,
\[
 B_k(0)=\tfrac12 c_k\log q+O_k(1),\qquad
 c_k=\sum_{j=1}^k(-1)^{j-1}\binom{k}{j}\sqrt j
 =\frac1{2\sqrt\pi}\int_0^\infty(1-e^{-v})^k\frac{dv}{v^{3/2}}>0 .
\]
The gamma terms are bounded at each of the finitely many arguments, and
the logarithmic derivatives of $L$ are bounded uniformly in $q$ by their
absolutely convergent prime-power series. The integral follows by inserting
$\sqrt j=(2\sqrt\pi)^{-1}\int_0^\infty(1-e^{-jv})v^{-3/2}dv$
into the finite sum. Thus every fixed $k$ eventually gives $B_k(0)>2$
as the conductor grows, regardless of the actual order of vanishing.
This obstruction concerns a fixed depth of this particular test.
The substitution $v=y\log k$ in the integral gives $c_k\sim(\pi\log k)^{-1/2}$.
The factor $(1-k^{-y})^k$ tends to $\mathbf1_{y>1}$ almost everywhere.
On $y\ge1/2$ the other factor $y^{-3/2}$ is integrable. For
$0<v<1/2$ bound $(1-e^{-v})^k$ by $v^k$, and for
$1/2\le v\le (\log k)/2$ bound it by $\exp(-\sqrt k)$.
These bounds make the omitted part negligible after multiplication by
$\sqrt{\log k}$. The surviving integral is $\int_1^\infty y^{-3/2}dy=2$.

\paragraph{A certified case outside Platt's all-moduli search range.}
Let $q=1\,000\,004\,329$ and $\chi(n)=(n/q)$. For this prime conductor
the entire strip $0<\Re s<1$, $|\Im s|\le1/5$ is zero-free for
$L(s,\chi)$. The included Python checker gives the outward enclosures
\[
   0.94936 < B_{10}(0)<1.66669,\qquad
   1.78105 < 2K_{10}(1/25) < 1.78106 .
\]
The modulus was selected by a preliminary numerical search. The proof
uses the final enclosure alone. Platt checked central nonvanishing for
all primitive characters through modulus $2\,000\,000$
(\href{https://doi.org/10.1090/mcom/3077}{Math. Comp. 85 (2016)}).
Platt and Trudgian subsequently studied even quadratic conductors up to $10^9$
(\href{https://arxiv.org/abs/2610.29018}{October 2026 preprint}).
The character here is even and lies just above that bound. We make no
priority claim for this individual conductor.

For reproducibility, at $s=1/2+\sqrt j$ the program evaluates
\[
  \frac{\Lambda_\chi'}{\Lambda_\chi}(s)
  =\tfrac12\bigl(\log(q/\pi)+\psi_0((s+b)/2)\bigr)
    -\sum_n\chi(n)\Lambda(n)n^{-s},
\]
where $\psi_0=\Gamma'/\Gamma$. Trial division verifies primality.
Euler's criterion evaluates the character at primes. All prime powers
through $N=65\,536$ are included. With $\Psi(x)=\sum_{n\le x}\Lambda(n)$
we have $\Psi(2m)-\Psi(m)\le\log\binom{2m}{m}\le 2m\log2$.
Summing at dyadic points and rounding $x$ up to the next such point yields
$\Psi(x)\le4x\log2$. Partial summation bounds the absolute tail by
$4(\log 2)sN^{1-s}/(s-1)$. For $\psi_0$ the recurrence shifts its argument
up by 128, then Binet's formula gives
\[
  \log y-\frac{1}{2y}-\frac{1}{12y^2}
     \le\psi_0(y)\le\log y-\frac{1}{2y}\qquad (y>0).
\]
Indeed the remainder is
$-2\int_0^\infty u\,du/((u^2+y^2)(e^{2\pi u}-1))$.
The checker uses $333/106<\pi<355/113$ and the standard library's
correctly rounded decimal arithmetic at 38 digits with outward expansion
after each operation. The finite calculation and the criterion give one
individual case and a reproducible central zero-free interval.
For background on what is known about the family see Soundararajan,
\href{https://doi.org/10.2307/2261390}{Ann. of Math. 152 (2000)},
and the recent work on pair correlations and Chowla's conjecture,
\href{https://doi.org/10.1007/s00208-026-03383-y}{Math. Ann. (2026)}.

\section{The best positive bound from prescribed samples}

Fix distinct offsets $a_1,\ldots,a_k>1/2$ and a target ordinate $t$. Write
$H_i=a_i\Re (\Lambda'_\chi/\Lambda_\chi)(1/2+a_i+it)$.
The preceding Hadamard calculation expresses $H_i$ as the sum over
$u=\gamma-t$ of $a_i^2/(a_i^2+u^2)$. Each admissible function
\[
 F(u)=\sum_{i=1}^k c_i\frac{a_i^2}{a_i^2+u^2},\qquad
 c_i\in\mathbb R,\quad F(0)=1,\quad F(u)\ge0\quad(u\in\mathbb R)
\]
gives the upper bound $m_t\le\sum c_iH_i$. The positivity condition
allows coefficients of either sign. Both the optimal certificate for a
given character and its asymptotic limitation can be computed exactly.

\paragraph{The optimal certificate.}
Form the real symmetric matrix
\[
 G_{ij}=\frac{a_jH_i+a_iH_j}{a_i+a_j},\qquad
 \mathbf1=(1,\ldots,1)^{\mathsf T}.
\]
Then $G$ is positive definite and the smallest admissible bound equals
\[
 \min_F\sum_\gamma F(\gamma-t)
    =\frac{1}{\mathbf1^{\mathsf T}G^{-1}\mathbf1}.
\]
Thus $G-2\mathbf1\mathbf1^{\mathsf T}$ failing to be positive
semidefinite is sufficient for quadratic central nonvanishing when $t=0$.

\emph{Proof.}
Use $g_i(u)=a_i/(a_i-iu)$, so $g_i(0)=1$, and symmetrize the zero
counting measure by giving mass $1/2$ at both $u$ and $-u$ for each
$u=\gamma-t$. This preserves the integral of any even function. Against
this measure the Gram matrix of the $g_i$ is real. The elementary identity
\[
 \Re\bigl(g_i(u)\overline{g_j(u)}\bigr)
 =\frac{a_j}{a_i+a_j}\frac{a_i^2}{a_i^2+u^2}
  +\frac{a_i}{a_i+a_j}\frac{a_j^2}{a_j^2+u^2}
\]
gives its entries. The integrals converge and positive definiteness
follows because a nonzero linear combination has finitely many zeros,
the poles are distinct, and the measure has infinite support.

Every admissible $F$ can be written $|\sum v_i g_i|^2$.
Indeed after multiplying by $\prod (u^2+a_i^2)$ its numerator is a real
polynomial of degree at most $2k-2$, nonnegative on $\mathbb R$.
Pair its conjugate nonreal roots and halve its real root multiplicities
to write it as $|P(u)|^2$ with $\deg P\le k-1$. Partial fractions
of $P(u)/\prod(a_i-iu)$ give the representation.
Multiplication by a phase makes $\sum v_i=1$.
Cauchy--Schwarz in the $G$ inner product gives the lower bound
$(\mathbf1^{\mathsf T}G^{-1}\mathbf1)^{-1}$, attained by
$v=G^{-1}\mathbf1/(\mathbf1^{\mathsf T}G^{-1}\mathbf1)$.
This vector is real, so its squared modulus is even. Its denominator
has distinct factors $u^2+a_i^2$ and its numerator has smaller degree,
hence the resulting function belongs to the stated span.
Finally a real vector with $v^{\mathsf T}(G-2\mathbf1\mathbf1^{\mathsf T})v<0$
has $\mathbf1^{\mathsf T}v\ne0$. Normalization then gives a bound below~2.

\paragraph{A sharp obstruction even with conductor-dependent coefficients.}
At $t=0$ set $S=\sum_i a_i^{-1}$ and hold the sites fixed. Uniformly over
primitive nonprincipal characters of conductor $q$, the minimum equals
\[
       \frac{\log q}{2S}+O_{a_1,\ldots,a_k}(1) \qquad(q\longrightarrow\infty).
\]
To see this write $G=(\log q)C+O(1)$ entrywise with
$C_{ij}=a_i a_j/(a_i+a_j)$. This uses absolute convergence of the
prime sums at all the sample sites. The matrix $C$ is positive definite since it is the Gram matrix of $x\mapsto a_i e^{-a_i x}$ on $[0,\infty)$.
For $w=C^{-1}\mathbf1$ put $R(z)=\sum_jw_j a_j/(z+a_j)$.
The equations $a_iR(a_i)=1$ give the rational identity
\[
 1-zR(z)=\prod_j\frac{a_j-z}{a_j+z}.
\]
Both sides have a common denominator, numerator degree at most $k$, the
same $k$ distinct zeros, and value $1$ at zero. Differentiation at zero
gives $\mathbf1^{\mathsf T}C^{-1}\mathbf1=2S$. Inversion of the
fixed-size matrix now yields the asymptotic. This rules out any
fixed choice of sampling sites as a uniform certificate, even if the
coefficients are selected optimally for each individual character.

The function attaining the least coefficient of $\log q$ is given by
$|R(-iu)/(2S)|^2$, or explicitly by
\[
 F_*(u)=\frac{\sin^2(\sum_j\arctan(u/a_j))}{S^2u^2},
 \qquad F_*(0)=1.
\]
For $k=1$ it is the single kernel $a_1^2/(a_1^2+u^2)$. For $k=2$ it is
the product of the two single kernels, as follows at once by the
sine addition formula. For $a_j=\sqrt j$ the optimal coefficient in
front of $\log q$ is $1/(2\sum_{j\le k}j^{-1/2})\sim 1/(4\sqrt{k})$.
The binomial test gives $c_k/2\sim1/(2\sqrt{\pi\log k})$.
These compare the leading coefficients for fixed sites as $q$ grows. The error bound depends on $k$.
The extremal calculation is an instance of a rational reproducing
kernel. For related classical theory see M.~M.~Dzhrbashyan,
\href{https://www.mathnet.ru/eng/sm3761}{\emph{Biorthogonal systems of
rational functions and best approximation of the Cauchy kernel on the
real axis}}, Math. USSR-Sb. 24 (1974), 409--433.

\paragraph{Allowing the sample sites to depend on the character.}
Let $M_k(\chi)$ be the infimum of the central bound over all distinct
$a_j>1/2$ and all admissible $F$ at those sites. For each fixed $k$,
uniformly over primitive nonprincipal characters as $q\to\infty$,
\[
                    \frac{M_k(\chi)}{\log q}\longrightarrow\frac{1}{4k}.
\]
Thus the obstruction persists even if both the sample locations and
the signed coefficients are selected after inspecting the character.

Here are details addressing sites that collide, escape to infinity, or
approach the edge of absolute convergence. Any such $F$ has the form
$P(u)/\prod_j(1+u^2/a_j^2)$, with $P$ even and nonnegative,
$P(0)=1$, and $\deg P\le 2k-2$. For fixed $b>1/2$ we consequently have
\[
 F(u)\ge(2b)^{-2k}f(u),\qquad
 f(u)=\frac{P(u)}{(1+u^2/b^2)^k},
\]
since each denominator factor is at most $1+4u^2\le
4b^2(1+u^2/b^2)$. The fixed-pole extremum above gives
$\int_{\mathbb R} f\ge\pi b/k$ also with this repeated pole. Indeed,
replace the $k$ occurrences of $b$ by distinct sites tending to $b$,
keep $P$ fixed, and use dominated convergence in the inequality
$\int F\ge\pi/(\sum_j1/a_j)$. The latter equality for the least
integral follows from
$\int a_j^2/(a_j^2+u^2)\,du=\pi a_j$
and the minimizing $C$-quadratic form above.

On this one fixed finite-dimensional space,
\[
 \left|\sum_\gamma f(\gamma)-
        \frac{\log q}{2\pi}\int_{\mathbb R} f(u)\,du\right|
      \le A_{b,k}\int_{\mathbb R}|f(u)|\,du .
\]
For the basis $(b^2+u^2)^{-j}$, $1\le j\le k$, this follows by
differentiating at $v=b^2$ the identity
\[
 \sum_\gamma\frac{1}{v+\gamma^2}
   =\frac{1}{\sqrt v}\Re\frac{\Lambda_\chi'}{\Lambda_\chi}
          (1/2+\sqrt v)
   =\frac{\log q}{2\sqrt v}+O_b(1).
\]
All derivatives through order $k-1$ of the error are uniformly bounded
in a real neighborhood of $b^2$. The differentiated prime-power series
converge absolutely there, and parity gives two possible gamma terms. The zero sums differentiate termwise
by local domination. Differentiating $\int_{\mathbb R}(v+u^2)^{-1}du
=\pi/\sqrt v$ yields the corresponding real integrals. Equivalence of
norms on the fixed finite-dimensional space proves the stated estimate.

Since $f\ge0$, taking the infimum of the lower bound and then $q\to\infty$
gives $\liminf M_k(\chi)/\log q\ge(2b)^{-2k}b/(2k)$.
Letting $b\downarrow1/2$ gives $1/(4k)$. Conversely take fixed distinct
$a_j$ arbitrarily close to $1/2$ from above and apply the fixed-site
minimum, whose leading coefficient tends to $1/(4k)$. This proves the
matching upper bound. The limit is for $k$ fixed and makes no assertion
about the regime $k=k(q)\to\infty$.

\section{Further research questions}
\begin{enumerate}
\item How small can the set of selected primes be in a counterexample to McNew's conjecture? The four-prime result above and our explicit construction leave a wide interval. A minimal example may require a different pattern of deletions.
\item How large is the natural-density gain over Rankin's set? The exchange optimizes a particular family of digit patterns. Larger digit alphabets or several interacting distinguished coordinates could give larger improvements.
\item For polynomial ratios, determine the scale of $\beta_q-m_q$. Section~5 bounds it by $q^{-q+O(1)}$ up to a polynomial factor. A lower bound on a comparable scale would require a much more economical construction than the previously reported strict example.
\item For quadratic characters, can conductor-dependent choices of sampling sites and uniform estimates on their logarithmic derivatives force a central multiplicity bound below two? Section~7 gives the exact obstruction at fixed $k$, even for adaptive sampling sites. A positive resolution by this approach would require a growing number of samples.
\end{enumerate}
\begin{thebibliography}{9}
\bibitem{mcnew} N. McNew, \emph{On sets of integers which contain no three terms in geometric progression}, Math. Comp. 84 (2015), 2893--2910. \href{https://doi.org/10.1090/mcom/2979}{Published article}, \href{https://www.nathanmcnew.com/GPFsets.pdf}{author's text, Conjecture 7.2}.
\bibitem{asada} M. Asada, E. Fourakis, S. Manski, N. McNew, S. J. Miller and G. Moreland, \emph{Subsets of $\F_q[x]$ free of 3-term geometric progressions}, Finite Fields Appl. 44 (2017), 135--147. \href{https://doi.org/10.1016/j.ffa.2016.10.002}{Published article}, \href{https://www.nathanmcnew.com/fqxramsey.pdf}{author manuscript, Questions 6.1--6.5}.
\bibitem{rankin} R. A. Rankin, \emph{Sets of integers containing not more than a given number of terms in arithmetical progression}, Proc. Roy. Soc. Edinburgh Sect. A 65 (1961), 332--344. \href{https://doi.org/10.1017/S0080454100017726}{Publisher}.
\bibitem{behrend} F. A. Behrend, \emph{On sets of integers which contain no three terms in arithmetical progression}, Proc. Natl. Acad. Sci. USA 32 (1946), 331--332. \href{https://doi.org/10.1073/pnas.32.12.331}{Publisher}.
\bibitem{sutherland} A. Sutherland, \emph{Verification files for the shorter Lean proof of RH and of GRH for Dirichlet $L$-functions}, 26 November 2026. \href{https://math.mit.edu/~drew/shortgrh/}{Verification report}. Source and pinned locations are linked in Section~1.
\end{thebibliography}
\end{document}
